\chapter{Embedded Systems Tailoring Profile}
\ChapterLabel{EmbeddedSystemsProfile}
\index{domain tailoring profile!embedded systems}
\index{embedded systems}
\index{hardware--software co-design}

\section{How to Use This Appendix}
This appendix tailors the generic architecture template for systems in which
software interacts directly with dedicated hardware, constrained computing
resources, physical processes, or real-time events. It applies to bare-metal
firmware, real-time operating system (RTOS) products, embedded Linux systems, programmable
logic, hardware--software products, and other cyber-physical systems.

The profile augments the core template; it does not create a second
architecture specification. Project-specific requirements, interfaces,
decisions, diagrams, tests, and evidence shall remain in their authoritative
core chapters. When this profile applies, retain this appendix as the
authoritative tailoring guide, record the applicability decision in
\ChapterRef{DevelopmentPlan}, and use the additional concerns below when
creating and reviewing the core content. Remove this appendix from an
instantiated project only when an authorized human records that the profile is
not applicable.

\section{Applicability}
\label{Section:EmbeddedProfileApplicability}
Apply this profile when one or more of the following conditions is present:

\begin{itemize}
\item software executes on a dedicated processor, microcontroller, system on
chip, field-programmable device, or purpose-built computing module;
\item behavior depends on interrupts, deadlines, deterministic scheduling,
physical signals, sensors, actuators, or control loops;
\item CPU, memory, storage, bus bandwidth, energy, or thermal resources are
materially constrained;
\item the product requires hardware bring-up, boot firmware, board support,
device drivers, manufacturing provisioning, calibration, or field firmware
updates; or
\item safety, reliability, environmental, physical-security, certification,
or component-lifecycle concerns affect the architecture.
\end{itemize}

The applicability decision shall identify the embedded platform class, the
reason the profile applies, the hardware revisions in scope, and any other
domain profiles used concurrently. A project may combine this profile with a
future safety-critical, high-security, cloud, mobile, or machine-learning
profile. Each profile remains guidance; controlled project facts stay in the
core chapters.

\section{Core-Chapter Tailoring Map}
Use Table~\ref{Table:EmbeddedTailoringMap} to place each embedded concern in
one authoritative location.

\small
\begin{tabularx}{\linewidth}{L{4.2cm}Y}
\caption{Embedded-Systems Core-Chapter Tailoring Map\label{Table:EmbeddedTailoringMap}}\\
\toprule
\textbf{Core Chapter} & \textbf{Embedded-System Additions}\\
\midrule
\endfirsthead
\multicolumn{2}{l}{\small\itshape Embedded-Systems Core-Chapter Tailoring Map (continued)}\\
\toprule
\textbf{Core Chapter} & \textbf{Embedded-System Additions}\\
\midrule
\endhead
\bottomrule
\endlastfoot
Development Plan & Record the profile decision, hardware and firmware roles,
toolchain ownership, hardware availability, prototypes, certification gates,
manufacturing dependencies, and component-lifecycle risks.\\
System Context & Define the hardware--software boundary, physical environment,
sensors, actuators, operators, service tools, external buses, power sources,
and upstream or downstream systems.\\
User Stories and Use Cases & Include installation, startup, shutdown,
calibration, diagnostics, maintenance, update, recovery, degraded operation,
manufacturing, and service interactions where applicable.\\
Requirements & Add timing, resource, power, interface, boot, recovery, safety,
security, reliability, environmental, diagnostics, update, manufacturing, and
regulatory obligations.\\
Architecture Views & Document hardware--software allocation, firmware layers,
execution contexts, states, timing, memory, deployment, data and signal flow,
power modes, faults, and recovery.\\
Implementation & Identify processor and board revisions, boot chain, BSP, HAL,
RTOS or operating system, compiler, linker, SDK, libraries, generated code,
linker scripts, build variants, image formats, and programming tools.\\
Tests, Verification, and Validation & Add target execution, static analysis,
timing analysis, SIL, PIL, HIL, fault injection, power, environmental,
hardware-revision, update, recovery, and retained-evidence activities.\\
Release & Identify released hardware revisions, firmware images, hashes,
configuration data, programming instructions, compatibility, rollback,
manufacturing release, service provisions, and known limitations.\\
\end{tabularx}
\normalsize

\section{Hardware--Software Boundary and Allocation}
The System Context and architecture chapters shall identify every processor,
memory device, programmable component, peripheral, sensor, actuator, bus,
power domain, and external physical interface that materially constrains the
software. Define which behavior is implemented in hardware, firmware,
programmable logic, operating-system services, application software, or an
external system. Record alternatives and trade-offs for significant
allocation decisions.

For each allocated function, identify its driving requirements, execution
platform, inputs, outputs, timing needs, failure behavior, update boundary,
and verification method. Do not treat a vendor module, BSP, HAL, RTOS, device
driver, or generated component as an unexamined implementation detail when it
affects required behavior.

\section{Embedded Requirement Categories}
Table~\ref{Table:EmbeddedRequirementCategories} is a checklist of common
embedded requirement categories. It does not introduce new requirement ID
families. Classify each requirement in the existing family that best describes
its primary obligation. Propose a new family only through the controlled
procedure in \ChapterRef{Requirements}. Apply the domain-profile translation
rule in Section~\ref{Section:DomainProfileRequirements}: every applicable
embedded concern shall map to a project requirement, an existing requirement,
or an approved Not Applicable decision with rationale.

\small
\begin{tabularx}{\linewidth}{L{3.7cm}Y}
\caption{Embedded-System Requirement Categories\label{Table:EmbeddedRequirementCategories}}\\
\toprule
\textbf{Category} & \textbf{Questions the Requirements Shall Resolve}\\
\midrule
\endfirsthead
\multicolumn{2}{l}{\small\itshape Embedded-System Requirement Categories (continued)}\\
\toprule
\textbf{Category} & \textbf{Questions the Requirements Shall Resolve}\\
\midrule
\endhead
\bottomrule
\endlastfoot
Functional behavior & What observable control, sensing, computation,
communication, diagnostic, maintenance, and recovery behavior is required?\\
Timing and real time & What deadlines, periods, latency, jitter, sequencing,
startup, shutdown, and worst-case response limits apply?\\
Processing resources & What CPU utilization, execution-time, core-affinity,
interrupt-latency, and scheduling limits apply?\\
Memory and storage & What flash, RAM, stack, heap, nonvolatile storage,
retention, endurance, and memory-protection limits apply?\\
Power and thermal & What operating, sleep, peak, average, battery-life,
wake-up, brownout, thermal, and power-transition behavior is required?\\
Hardware interfaces & What electrical levels, pin assignments, connectors,
buses, protocols, rates, units, ranges, scaling, resolution, endianness, and
error behavior apply?\\
Boot and initialization & What boot sources, ordering, self-tests, startup
times, configuration checks, and failure responses are required?\\
Reset and recovery & What events cause reset, retry, failover, degraded mode,
safe state, rollback, or operator intervention?\\
Safety & What hazards, safety functions, safe states, diagnostic coverage,
independence, and integrity constraints apply?\\
Security & What root of trust, secure boot, image authentication, device
identity, key provisioning, access control, debug-port, anti-rollback, and
physical-attack provisions apply?\\
Reliability and availability & What service life, duty cycle, fault tolerance,
watchdog, data-integrity, availability, maintainability, and recovery targets
apply?\\
Environmental & What temperature, humidity, vibration, shock, altitude,
electromagnetic, contamination, ingress, and storage conditions apply?\\
Diagnostics and observability & What logs, counters, health indicators, fault
codes, traces, test points, telemetry, and service interfaces are required?\\
Calibration and configuration & What values are configurable or calibrated;
who controls them; how are range, integrity, compatibility, and provenance
maintained?\\
Firmware update & What transport, authentication, atomicity, compatibility,
interruption, rollback, recovery, audit, and field-service behavior applies?\\
Manufacturing and provisioning & What programming, serialization, key
injection, calibration, end-of-line testing, traceability, and acceptance data
is required?\\
Regulatory and certification & What standards, evidence, configuration
controls, independence, analysis, and approval obligations apply?\\
Component lifecycle & What approved components, alternates, lead-time,
obsolescence, silicon errata, vendor support, and redesign constraints apply?\\
\end{tabularx}
\normalsize

Each embedded requirement shall retain the normal priority, lifecycle status,
In Version value, use-case mapping, and test-case mapping. Quantitative limits
shall state units, applicable operating conditions, tolerance, and verification
method. Avoid terms such as ``fast,'' ``low power,'' or ``robust'' without a
measurable acceptance boundary.

\section{Required Architecture Views}
Select views according to risk and complexity, and explain any omitted view
that would normally apply. Create diagrams with TikZ and store them under the
owning core chapter.

\small
\begin{tabularx}{\linewidth}{L{4.0cm}Y}
\caption{Candidate Embedded-System Architecture Views\label{Table:EmbeddedArchitectureViews}}\\
\toprule
\textbf{View} & \textbf{Content}\\
\midrule
\endfirsthead
\multicolumn{2}{l}{\small\itshape Candidate Embedded-System Architecture Views (continued)}\\
\toprule
\textbf{View} & \textbf{Content}\\
\midrule
\endhead
\bottomrule
\endlastfoot
Hardware block view & Processing, memory, programmable logic, peripherals,
clock, reset, power, sensors, actuators, and external connections.\\
Hardware--software allocation & Required functions allocated to hardware,
firmware layers, programmable logic, or external systems.\\
Firmware layering & Bootloader, BSP, HAL, drivers, RTOS or kernel, middleware,
services, application, and diagnostic boundaries.\\
Boot and update sequence & Trust chain, initialization order, self-tests,
configuration loading, update, rollback, and recovery paths.\\
Execution-context view & Cores, processes, tasks, threads, interrupt service
routines, priorities, scheduling, synchronization, and shared resources.\\
State and mode view & Operational, maintenance, calibration, degraded, safe,
sleep, update, and fault states with permitted transitions.\\
Timing view & Periods, deadlines, latency, jitter, sequencing, bus transactions,
control-loop timing, and worst-case paths.\\
Memory view & Address spaces, partitions, executable regions, stacks, heaps,
DMA buffers, retained data, memory protection, and nonvolatile layout.\\
Data and signal flow & Units, scaling, sample rates, filtering, validity,
ownership, persistence, and error propagation.\\
Deployment view & Firmware images and configuration deployed to processors,
cores, boards, hardware revisions, and manufacturing or field tools.\\
Power and thermal view & Power domains, modes, transitions, wake sources,
consumers, thermal limits, throttling, and shutdown behavior.\\
Fault-containment view & Fault-detection boundaries, watchdogs, redundancy,
degraded operation, safe states, containment, and recovery.\\
\end{tabularx}
\normalsize

\section{Interfaces and Protocols}
For every hardware or communication interface, document the physical and
logical contract together. Include connector and pin identity, direction,
electrical characteristics, protocol and version, topology, addressing,
message or register definition, data type, units, range, resolution, scaling,
endianness, update rate, timing, initialization, ownership, concurrency,
timeouts, integrity checks, invalid values, error handling, compatibility, and
verification method where applicable.

An interface table or referenced interface-control document shall have a
controlled identifier and version. Each interface requirement and test shall
link to that authority. Do not infer an electrical limit, register behavior,
reserved bit, or protocol tolerance without controlled evidence.

\section{Resource Budgets}
Maintain resource budgets whenever architecture feasibility depends on a
finite embedded resource. Table~\ref{Table:EmbeddedResourceBudget} is a generic
pattern; create one row for each controlled allocation and reconcile component
allocations with the system total.

\small
\begin{tabularx}{\linewidth}{L{2.2cm}L{2.5cm}L{2.2cm}L{2.3cm}Y}
\caption{Generic Embedded Resource Budget\label{Table:EmbeddedResourceBudget}}\\
\toprule
\textbf{Resource} & \textbf{Owner} & \textbf{Allocation} & \textbf{Measured or Estimated} & \textbf{Requirement / Evidence}\\
\midrule
\endfirsthead
\multicolumn{5}{l}{\small\itshape Generic Embedded Resource Budget (continued)}\\
\toprule
\textbf{Resource} & \textbf{Owner} & \textbf{Allocation} & \textbf{Measured or Estimated} & \textbf{Requirement / Evidence}\\
\midrule
\endhead
\bottomrule
\endlastfoot
CPU time & [task or component] & [limit and units] & [value, method, conditions] & [requirement and test IDs]\\
Flash & [image or component] & [bytes or percent] & [map-file value] & [requirement and evidence]\\
RAM & [task or component] & [bytes or percent] & [static and peak value] & [requirement and evidence]\\
Stack & [task or context] & [bytes] & [worst observed or analyzed] & [requirement and evidence]\\
Bus bandwidth & [interface or flow] & [rate and utilization] & [value and conditions] & [requirement and test IDs]\\
Power or energy & [mode or component] & [power, current, or energy] & [value and conditions] & [requirement and test IDs]\\
Thermal margin & [component or mode] & [temperature or margin] & [value and conditions] & [requirement and evidence]\\
\end{tabularx}
\normalsize

Every budget shall state the measurement or estimation conditions, margin
policy, uncertainty, owner, update trigger, and response to an exceeded
allocation. Budget evidence shall distinguish analysis, simulation,
measurement, and worst-case guarantee.

\section{Safety, Reliability, and Fault Management}
Identify hazards, failure modes, fault sources, detection mechanisms,
containment boundaries, watchdogs, redundancy, degraded modes, safe states,
recovery behavior, diagnostic coverage, and residual risk where applicable.
Trace every credited mechanism to requirements and verification evidence.

Do not claim a safety integrity level, diagnostic coverage, failure rate,
independence, or risk acceptance without the applicable analysis and
authorized approval. When a dedicated safety profile or regulated process
applies, use it together with this embedded profile and record precedence for
any conflicting rule.

\section{Embedded Security}
Define the trust boundary and lifecycle from manufacturing through
decommissioning. Address device identity, root of trust, secure boot, image
authentication, signing and key custody, provisioning, credential storage,
debug and test access, privileged interfaces, update authorization,
anti-rollback, protected configuration, security logging, recovery, physical
attack assumptions, vulnerability response, and secure decommissioning as
applicable.

Security controls shall be mapped to threats, requirements, tests, retained
evidence, and operational responsibilities. Secrets and production keys shall
not be placed in this architecture document, source repository, AI prompt, or
test artifact.

\section{Toolchain and Configuration Control}
Record the controlled hardware revisions, processor variants, compiler and
linker versions, optimization and warning settings, linker scripts, BSP, HAL,
RTOS or operating-system version, SDK, libraries, code generators, programming
tools, static-analysis tools, build variants, configuration sources, and image
formats. Identify which outputs are generated, how they are reproduced, and
which tool qualifications or validations apply.

A released firmware image shall be traceable to source revision, configuration,
toolchain, dependencies, hardware compatibility, build procedure, image hash,
and approval evidence. Changes to silicon revision, board revision, compiler,
optimization, RTOS, BSP, HAL, linker script, or generated-code input shall
trigger an impact assessment.

\section{Embedded Verification and Validation}
Select test environments according to the requirement and the confidence
needed. The verification plan may include review, static analysis, host unit
tests, software-in-the-loop (SIL), processor-in-the-loop (PIL), target tests,
hardware-in-the-loop (HIL), fault injection, timing and worst-case execution
analysis, stress, endurance, power, thermal, environmental, electromagnetic,
hardware-revision compatibility, boot, update, rollback, and recovery testing.

Each result shall identify the hardware revision, firmware version, build
configuration, tool and fixture versions, calibration state, environmental
conditions, input data, procedure, expected result, actual result, retained
evidence, and reviewer. Simulation or host-test evidence shall not be
represented as target evidence. Validation shall use representative intended
operation and stakeholders or authorized representatives.

\section{AI Tailoring Instructions}
\index{AI agent!embedded systems tailoring}
When this profile is Applicable or TBD, an AI agent shall read this appendix
with the task-specific core chapters. Before editing, it shall ask for or
confirm:

\begin{itemize}
\item the platform class: bare metal, RTOS, embedded Linux, programmable logic,
or mixed;
\item processor, board, programmable-device, and hardware revisions in scope;
\item boot chain, operating system or RTOS, BSP, HAL, SDK, compiler, linker,
and relevant third-party components;
\item external interfaces, sensors, actuators, physical environment, timing,
resource, power, safety, security, regulatory, manufacturing, update, and
service constraints; and
\item available schematics, datasheets, interface-control documents, hazard
analysis, threat analysis, test evidence, and authorized approvals.
\end{itemize}

The agent shall add embedded concerns to the initialization gap list and place
project facts only in their authoritative core chapters. It shall not invent
pinouts, register behavior, electrical limits, memory sizes, execution times,
power measurements, hardware compatibility, safety classifications, security
claims, certification evidence, or test results. It shall report missing
controlled hardware evidence and every decision requiring authorized human
approval.

Use the following request for an embedded project:

\begin{verbatim}
Read tgmAIAgentGuide.tex and tgmEmbeddedSystemsProfile.tex. Apply the
Embedded Systems Tailoring Profile to [Project Name]. Ask me to
confirm the project namespace, profile applicability, platform class,
hardware revisions, toolchain, interfaces, constraints, applicable
standards, and available evidence. Add embedded concerns to the gap
list before editing. After human review, place project-specific
requirements, architecture, diagrams, tests, and evidence only in
their authoritative core chapters. Do not invent hardware facts or
verification results. Build the complete document, inspect the log,
and visually inspect the affected PDF pages.
\end{verbatim}

\section{Embedded-Profile Completion Checklist}
Before completing embedded tailoring, confirm that:

\begin{itemize}
\item profile applicability and rationale are recorded;
\item the hardware--software boundary and function allocation are explicit;
\item timing, resources, power, interfaces, boot, recovery, safety, security,
environmental, update, manufacturing, and lifecycle concerns were assessed;
\item required embedded views and resource budgets exist or have documented
omission rationales;
\item hardware revisions, firmware versions, toolchain, build variants, and
released images are configuration controlled;
\item requirements, use cases, tests, traceability, verification, validation,
and release records have reciprocal links; and
\item the build succeeds and affected tables, diagrams, links, and page breaks
have been visually inspected.
\end{itemize}
